\begin{Resumen}
Este Proyecto Final de Ingeniería diseñó y desarrolló un prototipo de apoyo al análisis estructural de resonancias magnéticas lumbares. El sistema integra procesamiento de imágenes, segmentación anatómica asistida por inteligencia artificial en planos sagital y axial, mediciones geométricas, visualización superpuesta y una salida estructurada editable para revisión profesional. La evaluación sagital sobre la partición retenida de SPIDER obtuvo un Dice macro de 0,893 y superó el umbral definido por el proyecto. El módulo axial, evaluado sobre Sudirman y Al-Kafri, obtuvo un Dice macro oficial de 0,679 y no alcanzó ese umbral; por ello se conserva como candidato complementario y no como baseline consolidado. El producto también integra flujos de revisión, persistencia y exportación, con pruebas técnicas documentadas y dos sesiones de revisión profesional cualitativa. Estas evidencias no equivalen a validación clínica: no se compararon mediciones ni resultados diagnósticos del sistema contra profesionales sobre estudios comunes. Se concluye que el prototipo demuestra factibilidad técnica dentro de un alcance académico acotado, mientras que su exactitud clínica, desempeño en práctica real y viabilidad comercial observada permanecen sin demostrar.

\vspace{0.5em}
\noindent\textbf{Palabras clave:}\enspace resonancia magnética lumbar; segmentación médica; inteligencia artificial; mediciones anatómicas; revisión profesional.
\end{Resumen}
